\section*{Chapter \ref{ch:b1:unemployment-flows} --- Unemployment: Stocks and Flows}
% Numbers recomputed in code/bachelor-1/22-unemployment-flows/exercises/solutions.py (tested).

\begin{solution}{exo:b1:unemployment-flows:ustar}
$u^* = 0.02/(0.02 + 0.28) = 6.67\%$.
\end{solution}

\begin{solution}{exo:b1:unemployment-flows:dur}
$1/0.25 = 4$ months.
\end{solution}

\begin{solution}{exo:b1:unemployment-flows:read}
In 1933, at the bottom of the Depression, firms hired almost no one: a
collapse of $f$, after a wave of separations. In 2009 both moved, but by the
evidence of \cref{sec:b1:unemployment-flows:flows}, mostly $f$. In 1944 war
production and the armed forces absorbed almost every worker: vacancies
everywhere and a very high $f$.
\end{solution}

\begin{solution}{exo:b1:unemployment-flows:predict}
$0.02/0.22 = 9.09\%$.
\end{solution}

\begin{solution}{exo:b1:unemployment-flows:half}
$\ln 2/0.30 = 2.31$ months.
\end{solution}

\begin{solution}{exo:b1:unemployment-flows:which}
$u^*_A = 0.015/0.315 = 4.76\%$ and $u^*_B = 0.017/0.237 = 7.17\%$. The change in $s$ alone gives
$0.017/0.317 = 5.36\%$; the change in $f$ alone gives $0.015/0.235 = 6.38\%$. The fall in job
finding moved unemployment more: $1.62$ points against $0.6$.
\end{solution}

\begin{solution}{exo:b1:unemployment-flows:nsw}
$w = 1 + 0.4 + (0.4/0.2)(0.004 + 0.02 + 0.28) = 1.4 + 2 \times 0.304 = 2.008$.
\end{solution}

\begin{solution}{exo:b1:unemployment-flows:firing}
$u^* = 0.016/(0.016 + 0.224) = 6.67\%$, unchanged; the \omterm{def:b1:solow-model:speed}{half-life} rises from $2.31$ to
$\ln 2/0.24 = 2.89$ months. Spells are longer and jobs last longer.
\end{solution}

\begin{solution}{exo:b1:unemployment-flows:youth}
Germany $1.9$, Spain $2.35$, Sweden $2.87$, Poland $4.07$. The young separate
more often --- first jobs, temporary contracts, moves between education
and work --- so their $s$ is higher; their $f$ need not be lower.
\end{solution}

\begin{solution}{exo:b1:unemployment-flows:solve}
$\dot u = -(s + f)(u - u^*)$, so $u(t) - u^* = (u(0) - u^*)e^{-(s + f)t}$, which tends to zero for any $u(0)$; $u$
stays in $[0, 1]$ because $\dot u > 0$ at $u = 0$ and $\dot u < 0$ at $u = 1$.
\end{solution}

\begin{solution}{exo:b1:unemployment-flows:three}
In the \omterm{def:b1:solow-model:steady-state}{steady state} $sE = fU$ and $xU = yN$, so $U = (s/f)E$ and $N = (x/y)(s/f)E$; with
$E + U + N = 1$, $E = 1/[1 + s/f + (x/y)(s/f)]$. The unemployment rate $U/(E + U) = s/(s + f)$ is
unchanged: with these flows, inactivity changes the size of the labour
force, not the rate.
\end{solution}

\begin{solution}{exo:b1:unemployment-flows:refute}
Recessions in which the inflow into unemployment rises sharply while the
\omterm{def:b1:unemployment-flows:rates}{job-finding rate} barely moves; or a decomposition, like that of
\cref{exo:b1:unemployment-flows:which}, giving most of the rise in $u^*$ to $s$
across many recessions.
\end{solution}

\begin{solution}{pb:b1:unemployment-flows:brenland}
\textbf{1.} $4645 - 4366 = 279$ thousand, $279/4645 = 6.006\%$.

\textbf{2.} $s = uf/(1 - u) = 0.06006 \times 0.30/0.93994 = 0.01917$.

\textbf{3.} Inflow $0.01917 \times 4366 = 83.7$ thousand; outflow $0.30 \times 279 = 83.7$ thousand.

\textbf{4.} $1/0.30 = 3.33$ months.

\textbf{5.} $s + f = 0.31917$; \omterm{def:b1:solow-model:speed}{half-life} $\ln 2/0.31917 = 2.17$ months.

\textbf{6.} $0.02 \times 4366 = 87.3$ thousand.

\textbf{7.} $(279 + 87.3)/4645 = 7.886\%$.

\textbf{8.} $u(t) = 6.006\% + (7.886\% - 6.006\%)e^{-0.31917t}$.

\textbf{9.} $7.373\%$, $6.728\%$, $6.283\%$ and $6.047\%$.

\textbf{10.} Three in ten of the unemployed find a job each month: the pool
turns over fast.

\textbf{11.} A fall in $f$ --- if the laid-off workers' skills do not fit the
jobs on offer, or if the closure is part of a recession --- or duration
dependence among the laid-off.

\textbf{12.} $w = 1.2 + 0.5 + (0.5/0.1)(0.005 + 0.01917 + 0.30) = 3.3209$ thousand dalens a month.

\textbf{13.} $1.6209$, which equals $(e/q)(r + s + f)$.

\textbf{14.} It rises by $0.3$ to $3.6209$.

\textbf{15.} It falls by $5 \times 0.05 = 0.25$ to $3.0709$.

\textbf{16.} At a lower wage, workers would shirk: the queue does not make a
lower wage an incentive.

\textbf{17.} It lowers the rent $e/q$ needed, so the wage, and raises employment.

\textbf{18.} $0.01917/(0.01917 + 0.25) = 7.122\%$.

\textbf{19.} $0.07122 \times 4645 = 330.8$ thousand.

\textbf{20.} $4$ months.

\textbf{21.} $\ln 2/0.26917 = 2.58$ months.

\textbf{22.} If part of the fall in $f$ comes from cash in hand, it reflects
the relief of a borrowing constraint rather than a distorted incentive, so
the efficiency cost of the reform is smaller than the rise in $u^*$ alone
suggests. The same study finds that the longer searches do not lead to
better jobs (Card, Chetty and Weber 2007): the gain is not in match
quality. % ledger: b1:card2007-match

\textbf{23.} Compare the effect of extra benefits with that of an equal
lump sum that does not depend on staying unemployed, as in the Austrian
study, and the quality of the jobs found.

\textbf{24.} On both: lower separations and slower hiring.

\textbf{25.} The reform raises the steady-state rate from $6.0\%$ to $7.1\%$ (from
$6.006\%$ to $7.122\%$).
\end{solution}
